\documentclass{article}
\usepackage{amsmath,amsthm,amssymb}
\usepackage[margin=1in]{geometry}

\newcommand{\Cat}{\mathcal C}
\newcommand{\dinv}{\operatorname{dinv}}

\title{Rational \(q,t\)-Catalan Formula}
\author{}
\date{}

\begin{document}
\maketitle

\section{Purpose}

This item records the conjectural formula and finite checker from Graham
Hawkes's paper \emph{A conjectured formula for the rational
\(q,t\)-Catalan polynomial} (Annals of Combinatorics 28, 749--795, 2024;
arXiv:2208.00577).

The source script is
\[
  \texttt{Conjectures-and-Computations/qt-catalan/qt-conjecture.py}.
\]
The curated script
\[
  \texttt{code/check\_rational\_qt\_catalan\_formula.py}
\]
ports the same finite check into a command-line tool with compact output.

\section{Step Coordinates}

The checker uses the source script's step-coordinate convention.  For a
coprime pair \((r,n)\), put
\[
  \ell=r-1.
\]
A generated path is a sequence
\[
  x=(x_0,\ldots,x_{\ell-1})
\]
constructed from left to right subject to the rational Dyck bound
\[
  \sum_{a=0}^{i} x_a
  \le
  \left\lfloor\frac{(i+1)n}{\ell+1}\right\rfloor
  \qquad(0\le i<\ell).
\]
The script computes a degree statistic incrementally.  For
\[
  \beta(i,j;x)=\sum_{a=i}^{j}x_a-\frac{n(j-i+1)}{\ell+1},
\]
define
\[
  \gamma(i,j;x)=
  \begin{cases}
    \min(x_{i-1},\lfloor-\beta(i,j;x)\rfloor), & \beta(i,j;x)<0,\\
    \min(x_i,\lfloor\beta(i,j;x)\rfloor), & \beta(i,j;x)>0,\\
    0, & \beta(i,j;x)=0.
  \end{cases}
\]
When a new final entry is appended, the degree is increased by the corresponding
\(\gamma(k,j;x)\) terms with the new endpoint \(j\).  This is the source
script's statistic used in the monomial comparison.

The maximum area parameter is
\[
  M=\sum_{i=0}^{\ell-1}
  \left\lfloor\frac{(i+1)n}{\ell+1}\right\rfloor .
\]
For a generated path \(x\), the script uses
\[
  \operatorname{area}(x)
  =
  M-\sum_{p=0}^{\ell-1}\sum_{a=0}^{p}x_a.
\]

\section{Conjectural Formula Check}

The formula is expressed as a monomial-string identity.  The source script
first finds the unique closest point to the diagonal used by the conjecture:
an index \(q\) such that
\[
  (q+1)n\equiv 1 \pmod{\ell+1},
\]
recorded as
\[
  \left(q,\left\lfloor\frac{(q+1)n}{\ell+1}\right\rfloor\right).
\]
The distinguished subfamily \(T\) consists of the generated paths satisfying
\[
  \sum_{a=0}^{q}x_a
  =
  \left\lfloor\frac{(q+1)n}{\ell+1}\right\rfloor .
\]

Every generated path contributes one left-side monomial record
\[
  (\operatorname{area}(x),\,M-d(x)),
\]
where \(d(x)\) is the generated degree statistic.  Each path in \(T\) contributes
a monomial string to the right side.  If \(a=\operatorname{area}(x)\) and
\(d=d(x)\), then the positive interval is
\[
  (j,M-d),\qquad a\le j\le M-a-d,
\]
when \(a\le M-a-d\).  If \(M-a-d<a\), the script records the negative correction
interval
\[
  (j,M-d),\qquad M-a-d+1\le j<a.
\]
The finite check is the exact multiset identity
\[
  \text{positive interval multiset}
  =
  \text{all path monomials}+\text{negative correction multiset}.
\]

\section{Reproducible Runs}

The source script's built-in example is \((r,n)=(7,12)\).  Running
\[
  \texttt{python code/check\_rational\_qt\_catalan\_formula.py}
\]
checks that case.  In the curated run it generated 2,652 paths, found closest
point \((2,5)\), and verified
\[
  2666\text{ positive terms}
  =
  2652\text{ all-path terms}
  +
  14\text{ negative correction terms}.
\]
The status was \(\texttt{PASS}\).

The script also accepts multiple cases, for example
\[
  \texttt{python code/check\_rational\_qt\_catalan\_formula.py
  --case 3/5 --case 5/8 --case 7/12}.
\]
All listed cases must be coprime.  The source README explicitly warns that the
conjecture is not expected to hold for non-coprime parameters in general, and
the curated checker rejects non-coprime inputs.

\section{Status}

This item is conjectural.  It packages the finite checks for selected
relatively prime parameter pairs; it is separate from the computer-assisted
proof item for Lemma 2 and Lemma 3 of Section 9.

\end{document}
